\documentclass{article}
\usepackage{amsmath}
\usepackage{amssymb}
\usepackage{fullpage}
\usepackage{enumerate}
\usepackage{xcolor}
\usepackage{bm}
\pagestyle{empty}
\usepackage{graphicx}
\graphicspath{{C:\Users\steve\Pictures}}
\begin{document}

\noindent{{\Large \bf Fall 2025, 4100 Problem Set 3, Hayden Fuller}

\large

\bigskip
\noindent{Please complete all problems. For any proof questions, please structure your proof similar to those from lecture.

\bigskip

\begin{enumerate}

\item Given that $\textbf{b}_1$ and $\textbf{b}_2$ are two nonzero, linear independent vectors in $\mathbb{R}^3$ and that $T\in\mathcal{L}(\mathbb{R}^4,\mathbb{R}^3)$ is such that 
$$T\textbf{e}_1=\textbf{b}_1,~~T\textbf{e}_2=-\textbf{b}_1,~~T\textbf{e}_3=2\textbf{b}_1,~~{\rm and}~~T\textbf{e}_4=\textbf{b}_2$$
find the general solution to $T{\bf x}={\bf b}_1$
\\
\\$\bm x=(x_1,x_2,x_3,x_4)$
\\$T\bm x=x_+1\bm b_1+x_2(-\bm b_1)+x_3(2\bm b_1)+x_4\bm b_2=(x_1-x_2+2x_3)\bm b_1+x_4\bm b_2$
\\since $\bm b_1$,$\bm b_2$ are linearly independent, $T\bm x=\bm b_1$ iff $x_1-x_2+2x_3=1$, $x_4=0$
\\parameterize free variablesd $s=x_2$, $t=x_3$
\\$x_1=1+s-2t$, $x_4=0$
\\$\bm x=(1+s-2t, s, t, 0$, $s,t\in\mathbb{R}$
\\$\bm x=(1,0,0,0)+s(1,1,0,0)+t(-2,0,1,0)$, $s,t\in\mathbb{R}$
\\an affine 2-plane in $\mathbb{R}^4$

\newpage\item Let $V$ be the vector space $\mathbb{P}_3(\mathbb{R})$ endowed with the inner product, $\langle x^i,x^j\rangle=\begin{cases}
1~{\rm if}~i=j\\
0~{\rm if}~i\neq j
\end{cases}$ and let $S$ be a subspace of $V$ given by, $S=\{f\in\mathbb{P}_3: f(1)=f'(1)=0\}$. Find bases for $S$ and $S^{\perp}$. 
\\
\\$V=\mathbb{P}_3(\mathbb{R})$
\\monomial basis $\{1,x,x^2,x^3\}$
\\$f(x)=a+bx+cx^2+dx^3$
\\inner product makes monomial orthonormal so
\\$f,g$ with vectors $(a,b,c,d)$ and $(a\tilde a + b\tilde b + c\tilde c + d\tilde d)$
\\$\langle f,g\rangle = a\tilde a + b\tilde b + c\tilde c + d\tilde d$
\\$S=\{f\in\mathbb P_3: f(1)=f'(1)=0\}$ constraints:
\\$f(1)=a+b+c+d=0,\qquad f'(1)=b+2c+3d=0.$
\\From $b+2c+3d=0$ we get 
\\$b=-2c-3d$
\\$a=-(b+c+d)=c+2d$
\\coefficient vector: $(a,b,c,d)=(c+2d,\,-2c-3d,\,c,\,d)= c(1,-2,1,0)+d(2,-3,0,1).$
\\$\dim S=2$
\\a convenient basis of $S$ is $\,\;S=\operatorname{span}\{\,1-2x+x^2,\;\;2-3x+x^3\,\}$
\\$S^\perp$ is the orthogonal complement in $\mathbb R^4$
\\Let a general polynomial in $S^\perp$ have coefficients $(\alpha,\beta,\gamma,\delta)$
\\$\alpha-2\beta+\gamma=0,\qquad 2\alpha-3\beta+\delta=0$
\\$\beta=2\gamma-\delta,\qquad \alpha=3\gamma-2\delta$
\\$(\alpha,\beta,\gamma,\delta)=\gamma(3,2,1,0)+\delta(-2,-1,0,1)$
\\$\dim S^\perp=2$
\\basis is $S^\perp=\operatorname{span}\{\,3+2x+x^2,\;\;-2-x+x^3\,\}$

\bigskip

{\large\bf Bonus}: Find an integral formula for the above inner product

\newpage\item Let $A\in\mathbb{R}^{m\times n}$ and $B\in\mathbb{R}^{m\times l}$. Show that ${\rm range}(A)$ is orthogonal to ${\rm range}(B)$ if and only if $A^TB$ is the zero matrix.
\\
\\Let the columns of $A$ be $a_1,\dots,a_n\in\mathbb R^m$ and the columns of $B$ be $b_1,\dots,b_l\in\mathbb R^m$
\\${\rm range}(A)=\operatorname{span}\{a_1,\dots,a_n\}$
\\${\rm range}(B)=\operatorname{span}\{b_1,\dots,b_l\}$
\\$(A^TB)_{ij}=a_i^T b_j$
\\If ${\rm range}(A)$ is orthogonal to ${\rm range}(B)$, then every column $a_i$ is orthogonal to every column $b_j$
\\$a_i^T b_j=0$ for all $i,j$
\\every $A^TB$ is zero, so $A^TB=0$
\\if $A^TB=0$ then $a_i^T b_j=0$ for all $i,j$
\\any $x\in{\rm range}(A)$ and $y\in{\rm range}(B)$ are linear combinations $x=\sum_i \alpha_i a_i,\; y=\sum_j \beta_j b_j$
\\their inner product is $x^Ty=\sum_{i,j}\alpha_i\beta_j\,a_i^T b_j=0$
\\every vector in ${\rm range}(A)$ is orthogonal to every vector in ${\rm range}(B)$
\\${\rm range}(A)\perp{\rm range}(B)$

\newpage\item Compute the fundemental subspaces of the following matrix

$$A=\left[
\begin{array}{ccccc}
1 & 0 & 4 & -3 & 2\\
4 & -2  & 1 & -12 & 3\\
-3 & 2 & 3 & 9 & -1
\end{array}
\right]$$
\\
\\$\mathrm{rref}(A)=\begin{bmatrix}
1 & 0 & 4 & -3 & 2\\[4pt]
0 & 1 & 15/2 & 0 & 5/2\\[4pt]
0 & 0 & 0 & 0 & 0
\end{bmatrix}$
\\$\operatorname{rank}(A)=2$
\\1) Column space $\operatorname{range}(A)\subset\mathbb R^3$
\\Pivot columns are 1 and 2
\\baisis given by first two columns of $A$
\\$\operatorname{range}(A)=\operatorname{span}\{\,[1,\,4,\,-3]^T,\;[0,\,-2,\,2]^T$
\\$\dim\operatorname{range}(A)=2$
\\2) Row space $\operatorname{row}(A)\subset\mathbb R^5$
\\basis is two independent rows of $A$
\\$\operatorname{row}(A)=\operatorname{span}\{\,[1,\,0,\,4,\,-3,\,2],\;[4,\,-2,\,1,\,-12,\,3]$
\\$\dim\operatorname{row}(A)=2$
\\3) Nullspace $\mathcal N(A)=\{x\in\mathbb R^5:Ax=0\}$
\\From RREF free variables are $x_3,x_4,x_5$
\\basis $\mathcal N(A)=\operatorname{span}\left\{\begin{bmatrix}-4\\-15/2\\1\\0\\0\end{bmatrix}, \begin{bmatrix}3\\0\\0\\1\\0\end{bmatrix}, \begin{bmatrix}-2\\-5/2\\0\\0\\1\end{bmatrix}\right\}$
\\$\dim\mathcal N(A)=5- \operatorname{rank}(A)=3$
\\4) Left nullspace $\mathcal N(A^T)=\{y\in\mathbb R^3: A^Ty=0\}$
\\Find $A^T$ nullspace
\\$\mathcal N(A^T)=\operatorname{span}\{\;[-1,\,1,\,1]^T\;\}$
\\$\dim\mathcal N(A^T)=3-\operatorname{rank}(A)=1$

\newpage\item Let 
$$A= \left[
\begin{array}{ccccccc}
1 & -1 & 0 & 0 & 0 & 0 & 0\\
-1 & 0 & 0 & 1 & 0 &  0 & 0\\
0 & 1 & 0 & -1 & 0 &  0 & 0\\
0 & -1 & 0 & 0 & 0 & 0  & 1 \\
0  & 0  & -1  & 0  & 1  & 0  & 0 \\
0  & 0  & 1  & 0  &  0 & -1  & 0\\
0 & 0 & 0 & 0 & -1 & 1 & 0
\end{array} \right] $$

\begin{enumerate}[a.]
\item Draw a directed graph having $A$ as its incidence matrix
\\
\\$e_1:\; 2 \to 1$ (column 1 has $+1$ in row 1 and $-1$ in row 2)
\\$e_2:\; 1 \to 3$ (column 2: $-1$ at row 1, $+1$ at row 3)
\\$e_3:\; 5 \to 6$
\\$e_4:\; 3 \to 2$
\\$e_5:\; 7 \to 5$
\\$e_6:\; 6 \to 7$
\\$e_7:\; 7 \to 4$
\begin{verbatim}
  1
  ^\ 
  | \
e1|  e2
  |   v
  2 <- 3
   \_e4

   4
   ^
   |
  e7
   |
   7 --> 5
   ^     |
  e6     v
   6 <-- 
    \___e3
\end{verbatim}
\newpage\item Find bases for the four fundamental subspaces; ${\rm range}(A)$, ${\rm range}(A^T)$, ${\rm null}(A)$, and ${\rm null}(A^T)$
\\
\\$m$ vertices and $c$ connected components the rank is $m-c$
\\$m=7$, $c=2$, $\operatorname{rank}(A)=7-2=5$
\\$\dim\operatorname{range}(A)=5$
\\$\dim\operatorname{null}(A)=n-\operatorname{rank}=7-5=2$
\\$\dim\operatorname{null}(A^T)=m-\operatorname{rank}=2$
\\1) $\displaystyle \operatorname{null}(A)$ (edge-cycle space)
\\There are two independent cycles here:
\\$1\!-\!3\!-\!2\!-\!1$ uses edges $e_1,e_2,e_4$
\\Put $1$ on those three edge-positions and $0$ elsewhere:
\\$v_1 = [\,1,\;1,\;0,\;1,\;0,\;0,\;0\,]^T$
\\cycle on vertices $5\!-\!6\!-\!7\!-\!5$ uses edges $e_3,e_6,e_5$:
\\$v_2 = [\,0,\;0,\;1,\;0,\;1,\;1,\;0\,]^T$
\\these are linearly independent and span $\operatorname{null}(A)=\operatorname{span}\{\,v_1,v_2\,\}$
\\2) $\operatorname{null}(A^T)$ (vertex-potentials constant on components)
\\$A^T y=0$ means for every edge $j$, $y_{\text{head}(j)}-y_{\text{tail}(j)}=0$
\\$y$ is constant on each connected component
\\indicator of component $\{1,2,3\}$ is $u_1 = [\,1,\,1,\,1,\,0,\,0,\,0,\,0\,]^T$
\\indicator of component $\{4,5,6,7\}$ is $u_2 = [\,0,\,0,\,0,\,1,\,1,\,1,\,1\,]^T$
\\$\operatorname{null}(A^T)=\operatorname{span}\{\,u_1,u_2\,\}$
\\3) $\displaystyle \operatorname{range}(A)$ (column space)
\\$\operatorname{range}(A)$ is the orthogonal complement of $\operatorname{null}(A^T)$
\\consists of all vectors in $\mathbb R^7$ whose entries sum to zero on each connected component
\\basis for component $\{1,2,3\}$:  $r_1 = [\,1,-1,0,0,0,0,0\,]^T,\qquad r_2 = [\,0,1,-1,0,0,0,0\,]^T$
\\basis for component $\{4,5,6,7\}$: $r_3 = [\,0,0,0,1,-1,0,0\,]^T,\quad r_4 = [\,0,0,0,0,1,-1,0\,]^T,\quad r_5 = [\,0,0,0,0,0,1,-1\,]^T$
\\These 5 vectors are independent and span the column space $\operatorname{range}(A)=\operatorname{span}\{r_1,r_2,r_3,r_4,r_5\}$
\\4) $\displaystyle \operatorname{range}(A^T)$ (row space of $A$)
\\$\operatorname{range}(A^T)$ is the orthogonal complement of $\operatorname{null}(A)$
\\dimension 5
\\basis is given by any 5 independent rows of $A$
\\$\text{row}_1=[\,1,-1,0,0,0,0,0\,]$
\\$\text{row}_2=[\,-1,0,0,1,0,0,0\,]$
\\$\text{row}_3=[\,0,1,0,-1,0,0,0\,]$
\\$\text{row}_5=[\,0,0,-1,0,1,0,0\,]$
\\$\text{row}_6=[\,0,0,1,0,0,-1,0\,]$
\\$\operatorname{range}(A^T)=\operatorname{span}\{\text{row}_1,\dots,\text{row}_6\setminus\{\text{row}_4\}\}$
\\

\end{enumerate}


\newpage\item Let $T$ be the linear map on $\mathbb{R}^2$ defined as follows
\begin{eqnarray*}
T\left[
\begin{array}{cc}
1 \\
0
\end{array}\right]&=\left[
\begin{array}{cc}
-2 \\
1
\end{array}\right]\\
T\left[
\begin{array}{cc}
2 \\
1
\end{array}\right]&=\left[
\begin{array}{cc}
0 \\
4
\end{array}\right]
\end{eqnarray*}
\begin{enumerate}[a.]
\item Find $[T]_{\{\textbf{e}_i\}}$ the matrix representation of $T$ with respect to the standard basis.
\\
\\Let $A=[T]_{\{e_i\}}$. We are given
\\$A\begin{bmatrix}1\\0\end{bmatrix}=\begin{bmatrix}-2\\1\end{bmatrix}, \qquad A\begin{bmatrix}2\\1\end{bmatrix}=\begin{bmatrix}0\\4\end{bmatrix}$
\\the first column of $A$ is $[-2,1]^T$
\\columns of $A$ as $c_1,c_2$
\\second equation says $2c_1+c_2=[0,4]^T$
\\Using $c_1=[-2,1]^T$
\\$2[-2,1]^T + c_2 = [0,4]^T \implies c_2 = [0,4]^T - [-4,2]^T = [4,2]^T$
\\$[T]_{\{e_i\}} = A=\begin{bmatrix}-2 & 4\\[4pt] 1 & 2\end{bmatrix}$

\item Let $\{\textbf{v}_1,\textbf{v}_2\}$ be another basis for $\mathbb{R}^2$ where $\textbf{v}_1=\left[
\begin{array}{ccc}
1 \\
-1 
\end{array} \right]$ and 
$\textbf{v}_2=
\left[
\begin{array}{ccc}
2 \\
-3 
\end{array} \right]$. Find $[T]_{\{\textbf{v}_i\}}$, the matrix representation of $T$ with respect to the basis $\{\textbf{v}_i\}$ and use it write 
$T(4{\bf v}_1+5{\bf v}_2)$ as a linear combination of the basis vectors, $\textbf{v}_i$.
\\
\\$\mathbf v_1=[1,-1]^T,\ \mathbf v_2=[2,-3]^T$
\\Form $P=[\mathbf v_1\ \mathbf v_2]=\begin{bmatrix}1&2\\-1&-3\end{bmatrix}$
\\$[T]_{\{\mathbf v_i\}} = P^{-1}AP$
\\Compute $P^{-1}$. $\det P = 1\cdot(-3)-2\cdot(-1)=-1$
\\$P^{-1}=\frac{1}{-1}\begin{bmatrix}-3&-2\\[4pt]1&1\end{bmatrix}=\begin{bmatrix}3&2\\[4pt]-1&-1\end{bmatrix}$
\\$AP = \begin{bmatrix}-2&4\\[4pt]1&2\end{bmatrix}\begin{bmatrix}1&2\\-1&-3\end{bmatrix}= \begin{bmatrix}-6 & -16\\[4pt]-1 & -4\end{bmatrix}$
\\$[T]_{\{\mathbf v_i\}} = P^{-1}(AP)=\begin{bmatrix}3&2\\[4pt]-1&-1\end{bmatrix} \begin{bmatrix}-6 & -16\\[4pt]-1 & -4\end{bmatrix} =\begin{bmatrix}-20 & -56\\[4pt]7 & 20\end{bmatrix}$
\\$[T]_{\{\mathbf v_i\}}=\begin{bmatrix}-20 & -56\\[4pt]7 & 20\end{bmatrix}$
\\Apply $T$ to $4\mathbf v_1+5\mathbf v_2$
\\express in the $\{\mathbf v_i\}$-basis
\\coordinate vector of $w=4\mathbf v_1+5\mathbf v_2$ relative to $\{\mathbf v_i\}$ is $[4\; 5]^T$
\\Then the coordinate vector of $T(w)$ in that basis is
\\$[T]_{\{\mathbf v_i\}}\begin{bmatrix}4\\[4pt]5\end{bmatrix}
= \begin{bmatrix}-20 & -56\\[4pt]7 & 20\end{bmatrix}\begin{bmatrix}4\\[4pt]5\end{bmatrix}
= \begin{bmatrix}-20\cdot4-56\cdot5\\[4pt]7\cdot4+20\cdot5\end{bmatrix}
= \begin{bmatrix}-80-280\\[4pt]28+100\end{bmatrix}
= \begin{bmatrix}-360\\[4pt]128\end{bmatrix}$
\\$T(4\mathbf v_1+5\mathbf v_2) = -360\,\mathbf v_1 + 128\,\mathbf v_2$




\end{enumerate}


\newpage\item Let $T:\mathbb{P}_2\rightarrow\mathbb{R}_3$ be defined by $T(f)=f(-2){\bf e}_1+f(1){\bf e_2}+f(3){\bf e_3}$. Find the matrix representation of $T$ with respect to the standard bases of each space and use it to find the equation of the parabola passing through the points, $(-2,8),(1,3),(3,-5)$. 
\\
\\The standard basis of $\mathbb P_2$ is $\{1,x,x^2\}$
\\$f(-2)=a-2b+4c,\qquad f(1)=a+b+c,\qquad f(3)=a+3b+9c$
\\since $[f(-2),f(1),f(3)]^T=[T][a\; b\; c]^T$
\\find the parabola through $(-2,8),(1,3),(3,-5)$ solve $\begin{bmatrix}
1 & -2 & 4\\[4pt]
1 & 1 & 1\\[4pt]
1 & 3 & 9
\end{bmatrix}
\begin{bmatrix}a\\ b\\ c\end{bmatrix}
=
\begin{bmatrix}8\\3\\-5\end{bmatrix}$
\\$a=\frac{28}{5},\qquad b=-\frac{32}{15},\qquad c=-\frac{7}{15}$
\\$f(x)=\frac{28}{5}-\frac{32}{15}x-\frac{7}{15}x^2
=\frac{1}{15}\bigl(84-32x-7x^2\bigr)$







\end{enumerate}
\end{document}